% LaTeX companion of hw05.typ. Both formats are editable.
\section{Homework 5 - submitted work}

\subsection{Exercise 56 - linear independence}

Rearrange the four given vectors as

\[v_{1} = \begin{pmatrix}
e \\
1 \\
0 \\
0 \\
0 \\
0
\end{pmatrix},\quad v_{2} = \begin{pmatrix}
k \\
m \\
1 \\
0 \\
0 \\
0
\end{pmatrix},\quad v_{3} = \begin{pmatrix}
a \\
b \\
c \\
d \\
1 \\
0
\end{pmatrix},\quad v_{4} = \begin{pmatrix}
f \\
g \\
h \\
i \\
j \\
1
\end{pmatrix}.\]

Each \(v_{i}\) has an entry where all preceding vectors have \(0\) and it has a nonzero entry. Thus, in a relation \(c_{1}v_{1} + c_{2}v_{2} + c_{3}v_{3} + c_{4}v_{4} = 0\), the sixth coordinate forces \(c_{4} = 0\), and the same argument forces \(c_{3},c_{2},c_{1} = 0\) one by one. The vectors are linearly independent for all \(a,b,\ldots,m \in \mathbb{R}\).

\subsection{Exercise 33 - hyperplanes}

For \(c_{1}x_{1} + \ldots + c_{n}x_{n} = 0\), let \(A = \begin{pmatrix}
c_{1} & c_{2} & \ldots & c_{n}
\end{pmatrix}\). The hyperplane is \(\text{ker}\left( T_{A} \right)\), and the image of \(T_{A}:\mathbb{R}^{n}\rightarrow\mathbb{R}\) is nonzero because some \(c_{i} \neq 0\). Hence its image has dimension \(1\), and rank-nullity gives

\(\text{dim}(V) = n - 1.\)

Thus a hyperplane in \(\mathbb{R}^{3}\) is a plane, and a hyperplane in \(\mathbb{R}^{2}\) is a line.

\subsection{Exercise 63 - equal-dimensional nested subspaces}

Let \(\left( {v_{1},\ldots,v_{m}} \right)\) be a basis of \(V\), with \(V \subset W\) and \(\text{dim}(V) = \text{dim}(W) = m\). The vectors \(v_{1},\ldots,v_{m}\) lie in \(W\) and are linearly independent. By Theorem 3.3.4 they form a basis of \(W\). Every \(w \in W\) is therefore a linear combination of the \(v_{i}\), so \(w \in V\). Thus \(W \subset V\), and \(V = W\).

\subsection{Exercise 12 - arithmetic sequences}

Let \(S\) be the set of all arithmetic sequences. It contains the zero sequence. If

\(m = \left( {m_{0},m_{0} + k,m_{0} + 2k,\ldots} \right),\quad n = \left( {n_{0},n_{0} + l,n_{0} + 2l,\ldots} \right),\)

then

\(m + n = \left( {m_{0} + n_{0},m_{0} + n_{0} + \left( {k + l} \right),m_{0} + n_{0} + 2\left( {k + l} \right),\ldots} \right) \in S.\)

For \(r \in \mathbb{R}\),

\(r\left( {m_{0},m_{0} + k,m_{0} + 2k,\ldots} \right) = \left( {rm_{0},rm_{0} + rk,rm_{0} + 2rk,\ldots} \right) \in S.\)

Hence \(S\) is a subspace.

\subsection{\texorpdfstring{Exercise 28 - commuting \(2 \times 2\) matrices}{Exercise 28 - commuting 2 \textbackslash times 2 matrices}}

Let \(A = \begin{pmatrix}
a & b \\
c & d
\end{pmatrix}\) and \(B = \begin{pmatrix}
1 & 1 \\
0 & 1
\end{pmatrix}\). The equation \(AB = BA\) gives

\(\begin{pmatrix}
a & {a + b} \\
c & {c + d}
\end{pmatrix} = \begin{pmatrix}
{a + c} & {b + d} \\
c & d
\end{pmatrix},\)

so \(c = 0\) and \(a = d\). Therefore

\(S = \left\{ \begin{pmatrix}
a & b \\
0 & a
\end{pmatrix}\  \middle| \ a,b \in \mathbb{R} \right\} = \left\{ a\begin{pmatrix}
1 & 0 \\
0 & 1
\end{pmatrix} + b\begin{pmatrix}
0 & 1 \\
0 & 0
\end{pmatrix}\  \middle| \ a,b \in \mathbb{R} \right\}.\)

The displayed matrices are linearly independent, so they are a basis and \(\text{dim}(S) = 2\).

\section{Part A - Problem 6}

For each diagram, the submitted dependent triples are:

\textbf{(a) \(\left( {v_{1},v_{2},v_{3}} \right)\), \(\left( {v_{1},v_{2},v_{5}} \right)\), \(\left( {v_{1},v_{2},v_{4}} \right)\), \(\left( {v_{1},v_{3},v_{5}} \right)\), \(\left( {v_{1},v_{3},v_{4}} \right)\), \(\left( {v_{2},v_{3},v_{4}} \right)\), \(\left( {v_{2},v_{3},v_{5}} \right)\), \(\left( {v_{3},v_{4},v_{5}} \right)\), \(\left( {v_{2},v_{4},v_{5}} \right)\), \(\left( {v_{1},v_{4},v_{5}} \right)\) (10 sets).} (b) \(\left( {v_{1},v_{2},v_{3}} \right)\), \(\left( {v_{1},v_{2},v_{5}} \right)\), \(\left( {v_{1},v_{2},v_{4}} \right)\), \(\left( {v_{1},v_{3},v_{5}} \right)\), \(\left( {v_{2},v_{3},v_{4}} \right)\), \(\left( {v_{3},v_{4},v_{5}} \right)\), \(\left( {v_{2},v_{4},v_{5}} \right)\), \(\left( {v_{1},v_{4},v_{5}} \right)\) (8 sets). \textbf{(c) \(\left( {v_{1},v_{2},v_{3}} \right)\), \(\left( {v_{1},v_{2},v_{5}} \right)\), \(\left( {v_{1},v_{3},v_{5}} \right)\), \(\left( {v_{2},v_{3},v_{4}} \right)\), \(\left( {v_{3},v_{4},v_{5}} \right)\), \(\left( {v_{2},v_{4},v_{5}} \right)\) (6 sets).} (d) \(\left( {v_{1},v_{2},v_{3}} \right)\), \(\left( {v_{1},v_{3},v_{5}} \right)\), \(\left( {v_{1},v_{3},v_{4}} \right)\), \(\left( {v_{2},v_{3},v_{4}} \right)\), \(\left( {v_{2},v_{3},v_{5}} \right)\), \(\left( {v_{2},v_{4},v_{5}} \right)\).

\section{Part B}

\subsection{Problem 1 - images of independent lists}

\subsubsection{(a)}

False. Let \(T:\mathbb{R}^{2}\rightarrow\mathbb{R}^{2}\) be \(T(x) = \begin{pmatrix}
0 & 0 \\
0 & 0
\end{pmatrix}x\). The vectors \(\begin{pmatrix}
1 \\
0
\end{pmatrix},\begin{pmatrix}
0 \\
1
\end{pmatrix}\) are linearly independent, but their images are both \(0\), so the image list is linearly dependent.

\subsubsection{(b)}

True. Assume \(Y = \left( {T\left( x_{1} \right),\ldots,T\left( x_{k} \right)} \right)\) is linearly independent and \(d_{1}x_{1} + \ldots + d_{k}x_{k} = 0_{V}\). Then

\(d_{1}T\left( x_{1} \right) + \ldots + d_{k}T\left( x_{k} \right) = T\left( 0_{V} \right) = 0_{W}.\)

Independence of \(Y\) gives \(d_{1} = \ldots = d_{k} = 0\), so \(X\) is linearly independent.

\subsection{Problem 2 - prescribed kernel and image}

The rref required by

\(\text{ker}(T) = \left\{ x \in \mathbb{R}^{5}\  \middle| \ x_{1} = 5x_{2},x_{3} = 7x_{4} \right\}\)

is

\(\begin{pmatrix}
1 & {- 5} & 0 & 0 & 0 \\
0 & 0 & 1 & {- 7} & 0 \\
0 & 0 & 0 & 0 & 0
\end{pmatrix}.\)

The target image has basis \(\begin{pmatrix}
1 \\
0 \\
1
\end{pmatrix},\begin{pmatrix}
0 \\
1 \\
0
\end{pmatrix}\). Choosing the first and second nonredundant columns accordingly gives

\(A = \begin{pmatrix}
1 & {- 5} & 0 & 0 & 0 \\
0 & 0 & 1 & {- 7} & 0 \\
1 & {- 5} & 0 & 0 & 0
\end{pmatrix},\)

and \(T(x) = Ax\) is one solution.

The transformation is not unique. For example, elementary transformations preserve the rref, and

\(A' = \begin{pmatrix}
5 & {- 25} & 0 & 0 & 0 \\
0 & 0 & 1 & {- 7} & 0 \\
5 & {- 25} & 0 & 0 & 0
\end{pmatrix}\)

has the same required kernel and image but is different from \(A\).

\subsection{Problem 3 - maps defined on a basis}

\subsubsection{(a)}

Let \(\mathcal{B} = \left( {x_{1},\ldots,x_{n}} \right)\) be a basis of \(X\), and write \(v = c_{1}x_{1} + \ldots + c_{n}x_{n}\). Define

\(T(v) = c_{1}T\left( x_{1} \right) + \ldots + c_{n}T\left( x_{n} \right) = c_{1}y_{1} + \ldots + c_{n}y_{n}.\)

For \(v_{1} = \sum d_{i}x_{i}\), \(v_{2} = \sum e_{i}x_{i}\), this rule gives

\(T\left( {v_{1} + v_{2}} \right) = \sum\left( {d_{i} + e_{i}} \right)T\left( x_{i} \right) = T\left( v_{1} \right) + T\left( v_{2} \right)\)

and \(T\left( {kv_{1}} \right) = kT\left( v_{1} \right)\), so it is linear. If \(T'\) has the same values on the basis, then \(T'(v) = \sum c_{i}T'\left( x_{i} \right) = \sum c_{i}y_{i} = T(v)\), so it is unique.

\subsubsection{(b)}

Let \(\text{dim}(X) = n\) and let \(\left( {u_{1},\ldots,u_{k}} \right)\) be a basis of \(U\). Extend it to a basis \(\left( {u_{1},\ldots,u_{k},w_{k + 1},\ldots,w_{n}} \right)\) of \(X\). Since \(\text{dim}(V) = n - k\), choose a basis \(\left( {v_{1},\ldots,v_{n - k}} \right)\) of \(V\) and define

\(T_{U,V}\left( u_{i} \right) = 0,\quad T_{U,V}\left( w_{k + j} \right) = v_{j}.\)

By part (a), this is a valid linear transformation. Its image is \(V\) and its kernel is \(U\).

\subsubsection{(c)}

The map is not unique: the construction depends on an arbitrarily chosen basis of \(V\). Choosing a different basis can give different images for the \(w_{k + j}\) while retaining the required kernel and image.

\subsection{Problem 4 - ranks and nullities of a composition}

\subsubsection{(a)}

True. Since \(\text{im}\left( {S \circ T} \right) \subset \text{im}(S)\),

\(\text{rank}\left( {S \circ T} \right) = \text{dim}\left( {\text{im}\left( {S \circ T} \right)} \right) \leq \text{rank}(S).\)

\subsubsection{(b)}

True. View the composition in two stages. First \(T:U\rightarrow V\); then \(S:\text{im}(T)\rightarrow W\). Rank-nullity for the restricted second map gives

\(\text{rank}\left( {S \circ T} \right) = \text{rank}(T) - \text{dim}\left( {\text{ker}\left( S' \right)} \right) \leq \text{rank}(T).\)

\subsubsection{(c)}

True. Rank-nullity yields

\(\text{rank}(T) + \text{nullity}(T) = \text{rank}\left( {S \circ T} \right) + \text{nullity}\left( {S \circ T} \right).\)

Part (b) then implies \(\text{nullity}\left( {S \circ T} \right) \geq \text{nullity}(T)\).

\subsubsection{(d)}

False. If \(T\) is surjective, then \(S(V) = 0_{W}\), so \(\text{nullity}(S) = \text{dim}(V)\) and \(\text{nullity}\left( {S \circ T} \right) = \text{dim}(U)\). When \(\text{dim}(U) < \text{dim}(V)\), the claimed inequality fails.

\subsection{Problem 5 - symmetric and skew-symmetric matrices}

For \(T(A) = A + A^{T}\):

\subsubsection{(a)}

\(T\left( {A_{1} + A_{2}} \right) = \left( {A_{1} + A_{2}} \right) + \left( {A_{1} + A_{2}} \right)^{T} = T\left( A_{1} \right) + T\left( A_{2} \right)\), and \(T\left( {kA} \right) = kA + \left( {kA} \right)^{T} = kT(A)\), so \(T\) is linear.

\subsubsection{(b)}

If \(A \in \text{ker}(T)\), then \(A^{T} = - A\), so \(A \in \text{Skew}_{n}\). Conversely, \(B^{T} = - B\) implies \(T(B) = 0\), so \(\text{ker}(T) = \text{Skew}_{n}\). Also \(C = M + M^{T}\) satisfies \(C^{T} = C\), whence \(\text{im}(T) \subset \text{Sym}_{n}\). If \(D \in \text{Sym}_{n}\), choose \(N = \frac{1}{2}D\); then \(N + N^{T} = D\). Therefore \(\text{im}(T) = \text{Sym}_{n}\).

\subsubsection{(c)}

Both are subspaces. Each contains the zero matrix. If \(A^{T} = - A\) and \(B^{T} = - B\), then \(\left( {A + B} \right)^{T} = - \left( {A + B} \right)\) and \(\left( {kA} \right)^{T} = - kA\); this proves the subspace conditions for \(\text{Skew}_{n}\). Replacing \(- A, - B\) by \(A,B\) gives the same argument for \(\text{Sym}_{n}\).

\subsubsection{(d)}

There are \(n^{2}\) free entries in an arbitrary \(n \times n\) matrix. In a symmetric matrix, entries below the diagonal mirror entries above it, while the diagonal entries are free. Thus

\(\text{dim}\left( \text{Sym}_{n} \right) = \frac{n^{2} - n}{2} + n = \frac{n^{2} + n}{2}.\)

For a skew-symmetric matrix, the lower entries are the negations of the upper entries and every diagonal entry is zero. Thus

\(\text{dim}\left( \text{Skew}_{n} \right) = \frac{n^{2} - n}{2}.\)
